\begin{abstract}
\noindent
A language model can answer a medical question with a wrong dose as confidently as with a right one. We built a retrieval-augmented question answering system over 50 documents, 20 WHO guidelines and 30 open-access PubMed Central articles. It retrieves five passages with BGE-M3 embeddings fused with a from-scratch BM25, sends them to a free language model, and keeps the answer only when it cites the passages by number. It refuses with one fixed sentence when a cross-encoder judges that no passage answers the question, or when the model does not cite. We evaluate five systems, including two baselines (the model alone and BM25 retrieval), on 40 questions: 20 answerable, each with a verbatim evidence quote; 10 out of the corpus, used to choose the gates; and 10 near-domain questions written and committed before any run, as a held-out test. With no language model call, the MiniLM cross-encoder gate at its untuned 0.5 boundary answers all 20 answerable questions and refuses all 10 out-of-corpus ones (closest scores 0.870 and 0.156), where the cosine gate keeps a margin of 0.036. On the held-out questions, both gates let some questions through (3 of 10 for the cosine gate, 5 of 10 for the cross-encoder), and the model's own refusal caught them. The full evaluation is still running under a daily request quota; its numbers, the second model and the human ratings are marked as pending.
\end{abstract}
